\documentclass[14pt,a4paper]{extarticle}
\usepackage{fontspec}
\usepackage{polyglossia}
\setmainlanguage{russian}
\setmainfont{Times New Roman}
\setmonofont{Courier New}
\usepackage{setspace}
\onehalfspacing
\usepackage{geometry}
\geometry{left=30mm, right=15mm, top=20mm, bottom=20mm}
\usepackage{indentfirst}
\setlength{\parindent}{1.25cm}
\setlength{\parskip}{0pt}
\usepackage{xurl}
\urlstyle{same}
\usepackage{hyperref}
\usepackage{booktabs}
\usepackage{ragged2e}
\usepackage{graphicx}
\usepackage{float}
\usepackage{pdflscape}
\usepackage{listings}
\usepackage{xcolor}
\usepackage{array}
\usepackage{amsmath}
\usepackage{wasysym}
\allowdisplaybreaks[3]

\lstset{
language=Python,
basicstyle=\footnotesize,
keywordstyle=\color{blue},
stringstyle=\color{red},
commentstyle=\color{green!50!black},
showstringspaces=false,
breaklines=true,
frame=single,
numbers=left,
numberstyle=\tiny\color{gray},
stepnumber=1,
keepspaces=true,
extendedchars=true,
inputencoding=utf8
}

\newcommand{\num}[1]{#1}

\begin{document}
\sloppy
\thispagestyle{empty}
\centering{
Министерство образования и науки Российской Федерации\\
Федеральное государственное автономное образовательное учреждение высшего образования\\
«Санкт-Петербургский политехнический университет Петра Великого»\\[10pt]
Институт компьютерных наук и кибербезопасности\\
Высшая школа технологий искусственного интеллекта\\
02.03.01 Математика и компьютерные науки\\
\vspace{\fill}
{\large
Отчёт по дисциплине\\
{\Large \bfseries Вычислительная математика}\\
Лабораторная работа №4\\
{\Large <<Дискретное преобразование Фурье>>}\\
Вариант 18\\
\vspace{\fill}
}
{
\flushright
{\bfseries Выполнил:}\\
студент группы 5130201/40003\\
Джабраилов А.В.\\
\hspace{9cm}Подпись: \hrulefill \\
{\bfseries Принял:}\\
К. физ.-мат. наук, доц., доц.\\
Пак В.Г.\\
\hspace{9cm}Подпись: \hrulefill \\
\hspace{9cm}<<\rule{1.5cm}{0.4pt}>> \hrulefill\ 2026~г.
}
\centering{
Санкт-Петербург, 2026
}
}
\justifying
\newpage
\sloppy

\section*{1 Задание}
\begin{enumerate}
    \item Применяя дискретное преобразование Фурье для данного набора сигналов, получить вектор амплитуд, затем по нему обратным преобразованием получить исходный сигнал.
    \item Применяя быстрое преобразование Фурье с прореживанием по времени по основанию 2 для данного набора сигналов, получить векторы амплитуд по чётным и нечётным отсчётам, вектор амплитуд по всем отсчётам; затем по нему обратным преобразованием получить векторы обратных преобразований по чётным и нечётным отсчётам и исходный сигнал.
\end{enumerate}

\section*{2 Цель работы}
Освоить алгоритмы дискретного и быстрого преобразований Фурье, реализовать прямое и обратное преобразования для заданного дискретного сигнала и проверить восстановление исходных отсчётов.

\newpage

\section*{3 Ход работы}
Исходные данные варианта 18 содержат $n=16$ временных отсчётов.

\begin{table}[H]
\centering
\caption{Таблица исходных отсчётов.}
\small
\begin{tabular}{|c|c|c|c|c|c|c|c|c|}
\hline
$j$ & 0 & 1 & 2 & 3 & 4 & 5 & 6 & 7 \\ \hline
$y_j$ & $-0{,}78$ & $-1{,}22$ & $-1{,}34$ & $-1{,}39$ & $-1{,}42$ & $-1{,}43$ & $-1{,}42$ & $-1{,}41$ \\ \hline
\end{tabular}

\vspace{6pt}

\begin{tabular}{|c|c|c|c|c|c|c|c|c|}
\hline
$j$ & 8 & 9 & 10 & 11 & 12 & 13 & 14 & 15 \\ \hline
$y_j$ & $-1{,}37$ & $-1{,}30$ & $-1{,}10$ & $-0{,}10$ & $1{,}12$ & $1{,}25$ & $1{,}33$ & $1{,}36$ \\ \hline
\end{tabular}
\end{table}

Алгоритм вычислений:
\begin{enumerate}
    \item Вычислить вектор комплексных амплитуд $F_k$ по формуле прямого дискретного преобразования Фурье.
    \item По найденному вектору $F_k$ выполнить обратное дискретное преобразование Фурье и сравнить результат с исходным сигналом.
    \item Разделить исходный вектор на отсчёты с чётными и нечётными индексами.
    \item Для полученных подвекторов вычислить векторы амплитуд $F^0_k$ и $F^1_k$.
    \item По формулам быстрого преобразования Фурье получить общий вектор амплитуд $F_k$.
    \item Выполнить обратное БПФ по формуле (4), получить вспомогательные векторы $Y^0_j$, $Y^1_j$ и восстановить исходный сигнал.
\end{enumerate}

Для реализации использован язык \textit{Python}. В программе все преобразования вычисляются явно по формулам методических указаний, без использования готовой функции БПФ.

\subsection*{3.1 Дискретное преобразование Фурье}
Прямое дискретное преобразование Фурье определяется формулой:
\[
F_k=\sum_{j=0}^{n-1}y_j\exp\left(-\frac{2\pi i k}{n}j\right),
\quad k=0,1,\ldots,n-1.
\]
Обратное преобразование определяется формулой:
\[
y_j=\frac{1}{n}\sum_{k=0}^{n-1}F_k\exp\left(\frac{2\pi i k}{n}j\right),
\quad j=0,1,\ldots,n-1.
\]

Подставим $n=16$ и значения исходных отсчётов в формулу прямого ДПФ:
\begingroup\scriptsize
\begin{align*}
F_k={}&-0{,}78\exp\left(-\frac{2\pi i k}{16}\cdot0\right)
-1{,}22\exp\left(-\frac{2\pi i k}{16}\cdot1\right)
-1{,}34\exp\left(-\frac{2\pi i k}{16}\cdot2\right)\\
&-1{,}39\exp\left(-\frac{2\pi i k}{16}\cdot3\right)
-1{,}42\exp\left(-\frac{2\pi i k}{16}\cdot4\right)
-1{,}43\exp\left(-\frac{2\pi i k}{16}\cdot5\right)\\
&-1{,}42\exp\left(-\frac{2\pi i k}{16}\cdot6\right)
-1{,}41\exp\left(-\frac{2\pi i k}{16}\cdot7\right)
-1{,}37\exp\left(-\frac{2\pi i k}{16}\cdot8\right)\\
&-1{,}30\exp\left(-\frac{2\pi i k}{16}\cdot9\right)
-1{,}10\exp\left(-\frac{2\pi i k}{16}\cdot10\right)
-0{,}10\exp\left(-\frac{2\pi i k}{16}\cdot11\right)\\
&+1{,}12\exp\left(-\frac{2\pi i k}{16}\cdot12\right)
+1{,}25\exp\left(-\frac{2\pi i k}{16}\cdot13\right)\\
&+1{,}33\exp\left(-\frac{2\pi i k}{16}\cdot14\right)
+1{,}36\exp\left(-\frac{2\pi i k}{16}\cdot15\right).
\end{align*}
\endgroup

Например, для первых двух амплитуд:
\begingroup\small
\begin{align*}
F_0={}&-0{,}78-1{,}22-1{,}34-1{,}39-1{,}42-1{,}43-1{,}42-1{,}41\\
&-1{,}37-1{,}30-1{,}10-0{,}10+1{,}12+1{,}25+1{,}33+1{,}36=-9{,}220000,\\
F_1\approx{}&5{,}529825+9{,}351469i.
\end{align*}
\endgroup

Подставим найденные амплитуды в формулу обратного ДПФ:
\begingroup\scriptsize
\begin{align*}
\hat y_j=\frac{1}{16}\bigl(&-9{,}220000\exp\left(\frac{2\pi i\cdot0}{16}j\right)
+(5{,}529825+9{,}351469i)\exp\left(\frac{2\pi i\cdot1}{16}j\right)\\
&+(-2{,}486396+5{,}022864i)\exp\left(\frac{2\pi i\cdot2}{16}j\right)
+(-1{,}378383+0{,}540232i)\exp\left(\frac{2\pi i\cdot3}{16}j\right)\\
&+(0{,}080000+1{,}160000i)\exp\left(\frac{2\pi i\cdot4}{16}j\right)
+(-0{,}991293+1{,}391733i)\exp\left(\frac{2\pi i\cdot5}{16}j\right)\\
&+(-1{,}213604+0{,}322864i)\exp\left(\frac{2\pi i\cdot6}{16}j\right)
+(-0{,}800149+0{,}042971i)\exp\left(\frac{2\pi i\cdot7}{16}j\right)\\
&-0{,}740000\exp\left(\frac{2\pi i\cdot8}{16}j\right)
+(-0{,}800149-0{,}042971i)\exp\left(\frac{2\pi i\cdot9}{16}j\right)\\
&+(-1{,}213604-0{,}322864i)\exp\left(\frac{2\pi i\cdot10}{16}j\right)
+(-0{,}991293-1{,}391733i)\exp\left(\frac{2\pi i\cdot11}{16}j\right)\\
&+(0{,}080000-1{,}160000i)\exp\left(\frac{2\pi i\cdot12}{16}j\right)
+(-1{,}378383-0{,}540232i)\exp\left(\frac{2\pi i\cdot13}{16}j\right)\\
&+(-2{,}486396-5{,}022864i)\exp\left(\frac{2\pi i\cdot14}{16}j\right)
+(5{,}529825-9{,}351469i)\exp\left(\frac{2\pi i\cdot15}{16}j\right)\bigr).
\end{align*}
\endgroup

\subsection*{3.2 Быстрое преобразование Фурье}
Вектор отсчётов разделяется на два подвектора половинной длины:
\[
y^0_j=y_{2j},\qquad y^1_j=y_{2j+1},\quad j=0,1,\ldots,\frac{n}{2}-1.
\]
Для варианта 18:
\[
\bar y^0=(-0{,}78,\,-1{,}34,\,-1{,}42,\,-1{,}42,\,-1{,}37,\,-1{,}10,\,1{,}12,\,1{,}33),
\]
\[
\bar y^1=(-1{,}22,\,-1{,}39,\,-1{,}43,\,-1{,}41,\,-1{,}30,\,-0{,}10,\,1{,}25,\,1{,}36).
\]

Введём обозначение из методических указаний:
\[
W_n^{j\cdot k}=\exp\left(-\frac{2\pi i k}{n}j\right).
\]
Тогда для подвекторов:
\[
F^0_k=\sum_{j=0}^{\frac{n}{2}-1}y_{2j}W_{\frac{n}{2}}^{j\cdot k},
\qquad
F^1_k=\sum_{j=0}^{\frac{n}{2}-1}y_{2j+1}W_{\frac{n}{2}}^{j\cdot k}.
\]

Подставим $n=16$, то есть для подвекторов используется $W_8^{j\cdot k}$:
\begingroup\scriptsize
\begin{align*}
F^0_k={}&-0{,}78\exp\left(-\frac{2\pi i k}{8}\cdot0\right)
-1{,}34\exp\left(-\frac{2\pi i k}{8}\cdot1\right)
-1{,}42\exp\left(-\frac{2\pi i k}{8}\cdot2\right)\\
&-1{,}42\exp\left(-\frac{2\pi i k}{8}\cdot3\right)
-1{,}37\exp\left(-\frac{2\pi i k}{8}\cdot4\right)
-1{,}10\exp\left(-\frac{2\pi i k}{8}\cdot5\right)\\
&+1{,}12\exp\left(-\frac{2\pi i k}{8}\cdot6\right)
+1{,}33\exp\left(-\frac{2\pi i k}{8}\cdot7\right),\\
F^1_k={}&-1{,}22\exp\left(-\frac{2\pi i k}{8}\cdot0\right)
-1{,}39\exp\left(-\frac{2\pi i k}{8}\cdot1\right)
-1{,}43\exp\left(-\frac{2\pi i k}{8}\cdot2\right)\\
&-1{,}41\exp\left(-\frac{2\pi i k}{8}\cdot3\right)
-1{,}30\exp\left(-\frac{2\pi i k}{8}\cdot4\right)
-0{,}10\exp\left(-\frac{2\pi i k}{8}\cdot5\right)\\
&+1{,}25\exp\left(-\frac{2\pi i k}{8}\cdot6\right)
+1{,}36\exp\left(-\frac{2\pi i k}{8}\cdot7\right).
\end{align*}
\endgroup

По формулам БПФ:
\[
\begin{cases}
F_k=F^0_k+W_n^kF^1_k,\\
F_{k+\frac{n}{2}}=F^0_k-W_n^kF^1_k,
\end{cases}
\quad k=0,\ldots,\frac{n}{2}-1.
\]
Например, для $k=1$:
\begingroup\small
\begin{align*}
F_1&=(2{,}364838+4{,}654249i)
+\exp\left(-\frac{2\pi i}{16}\right)(1{,}126518+5{,}550854i)\\
&=5{,}529825+9{,}351469i,\\
F_9&=(2{,}364838+4{,}654249i)
-\exp\left(-\frac{2\pi i}{16}\right)(1{,}126518+5{,}550854i)\\
&=-0{,}800149-0{,}042971i.
\end{align*}
\endgroup

Обратное БПФ производится по формулам из методических указаний:
\[
\begin{cases}
y_j=Y^0_j+W_n^{-j}Y^1_j,\\
y_{j+\frac{n}{2}}=Y^0_j-W_n^{-j}Y^1_j,
\end{cases}
\quad j=0,\ldots,\frac{n}{2}-1,
\]
где
\[
Y^0_j=\frac{1}{n}\sum_{k=0}^{\frac{n}{2}-1}F^0_k W_{\frac{n}{2}}^{-j\cdot k},
\qquad
Y^1_j=\frac{1}{n}\sum_{k=0}^{\frac{n}{2}-1}F^1_k W_{\frac{n}{2}}^{-j\cdot k}.
\]

Для вычисления по этой формуле вектор амплитуд $F$ разделён на компоненты с чётными и нечётными индексами:
\begingroup\scriptsize
\[
\bar F^0=(F_0,F_2,F_4,F_6,F_8,F_{10},F_{12},F_{14}),
\quad
\bar F^1=(F_1,F_3,F_5,F_7,F_9,F_{11},F_{13},F_{15}).
\]
\endgroup

Подставим $n=16$. Например, для $j=0$:
\begingroup\scriptsize
\begin{align*}
Y^0_0={}&\frac{1}{16}\bigl((-9{,}220000)+(-2{,}486396+5{,}022864i)
+(0{,}080000+1{,}160000i)\\
&+(-1{,}213604+0{,}322864i)-0{,}740000
+(-1{,}213604-0{,}322864i)\\
&+(0{,}080000-1{,}160000i)+(-2{,}486396-5{,}022864i)\bigr)
=-1{,}075000,\\
Y^1_0={}&\frac{1}{16}\bigl((5{,}529825+9{,}351469i)
+(-1{,}378383+0{,}540232i)
+(-0{,}991293+1{,}391733i)\\
&+(-0{,}800149+0{,}042971i)
+(-0{,}800149-0{,}042971i)
+(-0{,}991293-1{,}391733i)\\
&+(-1{,}378383-0{,}540232i)
+(5{,}529825-9{,}351469i)\bigr)
=0{,}295000.
\end{align*}
\endgroup
Тогда
\[
y_0=Y^0_0+W_{16}^{0}Y^1_0=-1{,}075000+0{,}295000=-0{,}780000,
\]
\[
y_8=Y^0_0-W_{16}^{0}Y^1_0=-1{,}075000-0{,}295000=-1{,}370000.
\]

\newpage

\section*{4 Результаты}

\subsection*{4.1 Дискретное преобразование Фурье}
Вектор комплексных амплитуд, полученный по формуле прямого ДПФ:

\begin{table}[H]
\centering
\caption{Вектор амплитуд $F(k)$.}
\small
\begin{tabular}{|c|c||c|c|}
\hline
$k$ & $F(k)$ & $k$ & $F(k)$ \\ \hline
0 & $-9{,}220000$ & 8 & $-0{,}740000$ \\ \hline
1 & $5{,}529825+9{,}351469i$ & 9 & $-0{,}800149-0{,}042971i$ \\ \hline
2 & $-2{,}486396+5{,}022864i$ & 10 & $-1{,}213604-0{,}322864i$ \\ \hline
3 & $-1{,}378383+0{,}540232i$ & 11 & $-0{,}991293-1{,}391733i$ \\ \hline
4 & $0{,}080000+1{,}160000i$ & 12 & $0{,}080000-1{,}160000i$ \\ \hline
5 & $-0{,}991293+1{,}391733i$ & 13 & $-1{,}378383-0{,}540232i$ \\ \hline
6 & $-1{,}213604+0{,}322864i$ & 14 & $-2{,}486396-5{,}022864i$ \\ \hline
7 & $-0{,}800149+0{,}042971i$ & 15 & $5{,}529825-9{,}351469i$ \\ \hline
\end{tabular}
\end{table}

Результат обратного преобразования совпадает с исходным сигналом с точностью до машинной погрешности.

\begin{table}[H]
\centering
\caption{Проверка обратного ДПФ.}
\small
\begin{tabular}{|c|c|c|c|}
\hline
$j$ & $y_j$ & $\hat y_j$ & $|y_j-\hat y_j|$ \\ \hline
0 & $-0{,}780000$ & $-0{,}780000$ & $4{,}56\cdot10^{-15}$ \\ \hline
1 & $-1{,}220000$ & $-1{,}220000$ & $3{,}97\cdot10^{-15}$ \\ \hline
2 & $-1{,}340000$ & $-1{,}340000$ & $4{,}70\cdot10^{-15}$ \\ \hline
3 & $-1{,}390000$ & $-1{,}390000$ & $2{,}20\cdot10^{-15}$ \\ \hline
4 & $-1{,}420000$ & $-1{,}420000$ & $3{,}33\cdot10^{-15}$ \\ \hline
5 & $-1{,}430000$ & $-1{,}430000$ & $3{,}71\cdot10^{-15}$ \\ \hline
6 & $-1{,}420000$ & $-1{,}420000$ & $2{,}18\cdot10^{-15}$ \\ \hline
7 & $-1{,}410000$ & $-1{,}410000$ & $1{,}82\cdot10^{-15}$ \\ \hline
8 & $-1{,}370000$ & $-1{,}370000$ & $1{,}11\cdot10^{-15}$ \\ \hline
9 & $-1{,}300000$ & $-1{,}300000$ & $4{,}60\cdot10^{-15}$ \\ \hline
10 & $-1{,}100000$ & $-1{,}100000$ & $3{,}30\cdot10^{-15}$ \\ \hline
11 & $-0{,}100000$ & $-0{,}100000$ & $2{,}37\cdot10^{-15}$ \\ \hline
12 & $1{,}120000$ & $1{,}120000$ & $3{,}34\cdot10^{-15}$ \\ \hline
13 & $1{,}250000$ & $1{,}250000$ & $2{,}26\cdot10^{-15}$ \\ \hline
14 & $1{,}330000$ & $1{,}330000$ & $7{,}34\cdot10^{-15}$ \\ \hline
15 & $1{,}360000$ & $1{,}360000$ & $1{,}57\cdot10^{-15}$ \\ \hline
\end{tabular}
\end{table}

\subsection*{4.2 Быстрое преобразование Фурье}
После разделения исходного сигнала на отсчёты с чётными и нечётными индексами получены следующие вспомогательные векторы амплитуд:

\begin{table}[H]
\centering
\caption{Векторы амплитуд $F^0(k)$ и $F^1(k)$.}
\small
\begin{tabular}{|c|c|c|}
\hline
$k$ & $F^0(k)$ & $F^1(k)$ \\ \hline
0 & $-4{,}980000$ & $-4{,}240000$ \\ \hline
1 & $2{,}364838+4{,}654249i$ & $1{,}126518+5{,}550854i$ \\ \hline
2 & $-1{,}850000+2{,}350000i$ & $-2{,}340000+1{,}440000i$ \\ \hline
3 & $-1{,}184838-0{,}425751i$ & $-0{,}966518+0{,}190854i$ \\ \hline
4 & $0{,}080000$ & $-1{,}160000$ \\ \hline
5 & $-1{,}184838+0{,}425751i$ & $-0{,}966518-0{,}190854i$ \\ \hline
6 & $-1{,}850000-2{,}350000i$ & $-2{,}340000-1{,}440000i$ \\ \hline
7 & $2{,}364838-4{,}654249i$ & $1{,}126518-5{,}550854i$ \\ \hline
\end{tabular}
\end{table}

По формулам БПФ получен тот же вектор амплитуд $F(k)$, что и при прямом ДПФ (таблица 2). Максимальное отличие между прямым ДПФ и БПФ составило $6{,}58\cdot10^{-14}$.

Вспомогательные векторы обратного БПФ, вычисленные по формуле (4):
\begin{table}[H]
\centering
\caption{Векторы $Y^0(j)$ и $Y^1(j)$.}
\small
\begin{tabular}{|c|c|c|}
\hline
$j$ & $Y^0(j)$ & $Y^1(j)$ \\ \hline
0 & $-1{,}075000$ & $0{,}295000$ \\ \hline
1 & $-1{,}260000$ & $0{,}036955-0{,}015307i$ \\ \hline
2 & $-1{,}220000$ & $-0{,}084853+0{,}084853i$ \\ \hline
3 & $-0{,}745000$ & $-0{,}246831+0{,}595902i$ \\ \hline
4 & $-0{,}150000$ & $1{,}270000i$ \\ \hline
5 & $-0{,}090000$ & $0{,}512796+1{,}237999i$ \\ \hline
6 & $-0{,}045000$ & $0{,}972272+0{,}972272i$ \\ \hline
7 & $-0{,}025000$ & $1{,}279573+0{,}530017i$ \\ \hline
\end{tabular}
\end{table}

По формуле (4) восстановлен исходный сигнал. Максимальная погрешность восстановления составила $3{,}78\cdot10^{-15}$.

Как видно из таблиц, обратные преобразования возвращают исходные отсчёты, а отклонения имеют порядок машинной точности.

\newpage

\section*{5 Приложение}
\lstinputlisting[caption={Файл main.py}, label={lst:lab5code}]{../program/main.py}

\end{document}
